\section*{Anexo F. Análisis de robustez}
\phantomsection
\addcontentsline{toc}{subsection}{Anexo F. Análisis de robustez}
\setcounter{table}{0}
\renewcommand{\thetable}{F.\arabic{table}}
\renewcommand{\theHtable}{F.\arabic{table}}
\setcounter{figure}{0}
\renewcommand{\thefigure}{F.\arabic{figure}}
\renewcommand{\theHfigure}{F.\arabic{figure}}
\setcounter{equation}{0}
\renewcommand{\theequation}{F.\arabic{equation}}
\renewcommand{\theHequation}{F.\arabic{equation}}
% Entrega 3, fase 2 (23.10.2026; plan del 14.10, movimientos 1, 2, 4, 7 y 8): este anexo junta el analisis de
% robustez que estaba en capitulo5.tex (l. 1109-1399 del commit 7181ab8, seccion "Robustez de los resultados"),
% capitulo3.tex (3.7.7, sec:diseno-robustez), capitulo4.tex (4.7.5, sec:remedicion) y el bloque "Tablas del
% analisis de robustez" del Anexo A (anexos.tex l. 134-197), mas el barrido de demanda del corredor
% (capitulo5.tex l. 865-948). El texto que viene del Cap. V se movio sin reescribirlo: solo cambian las remisiones
% que ya no valen fuera del capitulo y se quitan los rotulos en negrita que repetian el titulo de su apartado. Los
% nombres del glosario se aplicaron en la pasada del glosario (23.10.2026, tarea 3f): IPPO con mensajes, sin mensajes
% y con penalizacion de spillback, fijo ajustado, fijo ajustado desde 10 s, ciclo minimo, medicion mas estricta,
% version evaluada y resumen en un numero. El material de los Caps. III y IV y del Anexo A se integro el
% 23.10.2026 sin reescribirlo: cada tramo lleva un comentario "Integracion" con su origen y con lo que se ajusto
% (remisiones, opciones de linea de ordenes dichas en palabras y frases que repetian lo ya dicho en este anexo).
% Ninguna etiqueta despues de los titulos con asterisco; tablas, figuras y ecuaciones llevan su etiqueta despues del
% caption y numeran F.n.
% Reduccion del Anexo F (23.10.2026; plan del 14.10, aprobado por Fabian: de la robustez por modelo queda solo la
% tabla de IPPO con mensajes): de las 18 tablas por modelo y de entrenamientos (F.6 a F.23) quedan tres; el detalle
% esta en el comentario de "Tablas por modelo". El resto del anexo no cambia, salvo la remision del parrafo del ciclo
% minimo en "Resumen en un numero y reparto de la variacion".

Este anexo reúne el análisis de robustez completo: qué cambia la
medición más estricta, el fijo ajustado desde 10~s de cada caso, la
ejecución, los veredictos de todos los modelos con esa medición, la
lectura por caso, el resumen en un número y el reparto de la variación,
las comprobaciones, lo que queda en pie, el barrido de demanda en el
corredor sin reentrenar y tres tablas por modelo y de entrenamientos;
las demás están en el repositorio del proyecto. Su resumen, con el
resultado de IPPO con mensajes en cada caso, está en la
Sección~\ref{sec:robustez}. El análisis se decidió después
del veredicto: P1 a P5 se congelaron el 25.09.2026, y P6 se cerró el
28.09.2026 con la misma regla.
% fuente: fechas: plan2/veredicto_congelado_0355.csv (P1-P5, 25.09) y plan2/veredicto.csv (P6, 28.09).
% Verificacion de las tareas 4 y 5 (23.10.2026): (1) decia "y las tablas por modelo"; tras la reduccion del anexo
% quedan tres (ciclo minimo, IPPO con mensajes y entrenamientos de Double DQN) y las demas estan en el repositorio
% ("Tablas por modelo"); (2) sale la frase "Se remite tambien a este anexo desde el diseno del analisis
% (sec:diseno-robustez), desde su ejecucion (sec:remedicion), desde el efecto de la demanda (sec:sensibilidad-demanda)
% y desde las conclusiones y las recomendaciones", que era una nota de edicion
% Pasada del glosario (23.10.2026, tarea 3f): la prosa de este anexo usa los nombres del glosario del plan del 14.10
% (seccion 3.5) y se quito la lista de nombres anteriores que iba aqui, porque las tablas se regeneran con los nombres
% nuevos. images/fig_robustez_mejora_caso.png se regenero el 23.10.2026 con los nombres nuevos. images/fig_plan2_demanda.png
% y tablas/demanda_corredor.tex se regeneraron despues, el mismo 23.10.2026, con barrido_demanda.py ya pasado a los
% nombres del glosario; se quito la frase del caption de fig:plan2-demanda que aclaraba los rotulos viejos

% Integracion (23.10.2026): aqui iban capitulo3.tex l. 1557-1563 (comienzo de 3.7.7, sec:diseno-robustez) y
% capitulo4.tex l. 1288-1297 (comienzo de 4.7.5, sec:remedicion) del commit 7181ab8. No se repiten: todo lo que
% dicen (tres decisiones del metodo revisadas despues del cierre, sin reentrenar; la metrica, el fijo barrido desde el
% verde minimo y el control trivial; que no reemplaza al veredicto; resultados en sec:robustez y ejecucion en
% sec:remedicion) ya esta en esta entrada, en "Que cambia" y en "Ejecucion". El texto sigue en el commit 7181ab8.

% Verificacion de la fase 2 (23.10.2026): la entrada y "Que cambia" se reescribieron para que cada idea vaya una vez
% (antes los tres cambios se contaban tres veces, que P6 no se vuelve a barrer cuatro, la decision tomada al ver P3
% dos y que las versiones evaluadas son las del plan dos). Los tres puntos que encontro la revision pasan a "Que
% cambia", cada uno con su decision; el parrafo "La medicion mas estricta cambia tres cosas" se quito, porque sus
% programas ya estan en "Fijo ajustado desde 10 s" y su baja del fijo se da alli una sola vez, con sus semillas. Se
% aplica tambien el glosario del plan del 14.10 (rejilla de valores por malla del barrido; P1 a P6 por los nombres
% internos de escenario; version guardada, entrenamiento, control y cambio frente al fijo).
Después de cerrar el plan, una revisión del alcance buscó de qué
dependían las cifras del veredicto y encontró tres puntos que la
especificación no cubría, que se describen más abajo, en ``Qué cambia''.
Este anexo repite el análisis de P1 a P5 con esos tres puntos
cambiados, sin reentrenar ningún modelo, y lo extiende a P6 una vez
cerrado su fijo. La especificación del plan pedía abrir la rejilla de
valores del barrido hasta que el óptimo quedara en su interior, con el
verde mínimo de 10~s como piso válido, y los ganadores del primer
barrido cumplían esa regla: en P3, por ejemplo, 25~s de avenida y 15~s
de transversal, dentro de una rejilla de 15 a 35~s y de 10 a 20~s.
Ampliar la rejilla hasta 10~s en la avenida se decidió después de ver
que la versión evaluada de IPPO con mensajes en P3 daba el mismo
resultado que el ciclo mínimo, y la especificación no lo exigía. Por
eso esta medición más estricta es un análisis de robustez: su
referencia se eligió después de ver los resultados, y no reemplaza el
veredicto pre-registrado.
% fuente: plan2/ESPEC.md, seccion 3 ("malla abierta hasta que el optimo sea interior (el verde minimo 10 s es un
% piso valido)"); comentario de plan2/red.add.xml en el commit 37f9680 (malla 15-35 x 10-20; ganador 25/15, ciclo
% 46 s, desfases 0/23/0)

\subsection*{Qué cambia}

% Integracion (23.10.2026): los cuatro parrafos siguientes son capitulo3.tex l. 1565-1598 del commit 7181ab8 (3.7.7),
% con su fuente (l. 1599-1606). Sin reescribir; solo cambian: (1) salen las dos frases que repetian la decision de
% ampliar la malla a 10 s despues de ver P3 ("Ampliar la malla hacia 10 s se decidio despues de ver que la politica
% publicada de IPPO vecinos en P3 era el programa trivial. La regla del plan no lo exigia, porque ningun ganador del
% barrido del plan habia quedado en el borde de su malla."), que ya esta en la entrada; (2) la remision a
% sec:remedicion de l. 1581 pasa a texto interno ("mas abajo, en Fijo ajustado desde 10 s", donde esta lo de P6).
% Verificacion de la fase 2 (23.10.2026): cada decision lleva aqui el punto de la revision que la motivo (antes en la
% entrada) y lo que antes decia "La medicion mas estricta cambia tres cosas" (vehiculos que no entran, semillas de
% ajuste, ciclo de 26 s); P6 se trata una sola vez, en "Fijo ajustado desde 10 s".
La primera decisión es la métrica. En P2 la espera para entrar a la
red, que el \texttt{timeLoss} no incluye, es mayor que el propio
\texttt{timeLoss}: con el fijo ajustado, 239.6 frente a 154.7~s/veh, y
los controles adaptativos bajan el \texttt{timeLoss} en parte con cola
que queda fuera de la red (Sección~\ref{sec:formulacion}). En
el análisis de robustez la regla de veredicto se aplica, por eso, al
retraso total (Ecuación~\ref{eq:retraso-total}), que suma al
\texttt{timeLoss} la espera para entrar a la red y cuenta también a los
vehículos que no llegan a entrar.
% fuente: plan2/evaluacion_corredor_alta_resumen.csv (fijo: retraso_insercion_prom 239.5657, timeloss_prom
% 154.7013)

La segunda es el fijo de referencia. El barrido del plan partía de
12~s de verde en la avenida en P1 y de 15~s en P2 a P5, por encima del
verde mínimo del agente, de 10~s. En el análisis de robustez se vuelve a
barrer con la rejilla de los dos verdes desde el mínimo de 10~s, con el
retraso total como criterio y con las mismas semillas de ajuste, 2001 a
2003; la regla de interioridad y el techo de 60~s no cambian. El
resultado, el fijo ajustado desde 10~s de cada caso, está más abajo, en
``Fijo ajustado desde 10~s''.
% fuente: mallas del 25.09: comentarios de plan2/<esc>.add.xml en el commit 37f9680 (corredor: avenida desde 12 s;
% corredor_alta, red, red_alta y malla3: avenida 15/20/25/30/35 s)

La tercera es un control nuevo, el ciclo mínimo
(Secciones~\ref{sec:modelos} y~\ref{sec:programa-trivial}): cada verde
en 10~s, los ámbares de 3~s y todos los desfases en cero, un ciclo de
26~s. El plan no tenía un control que dijera si un aprendiz gana solo
por cambiar de fase en cuanto la máscara lo permite. El ciclo mínimo se
evalúa en cada caso con las semillas 1001 a 1030 y se juzga frente al
fijo ajustado desde 10~s, como el actuado. La versión evaluada que da
la misma media de retraso total y la misma media de espera detenida que
el ciclo mínimo, con una diferencia de a lo sumo 0.0005~s/veh en cada
una, se marca como ``igual al ciclo mínimo''. La marca no cambia el
veredicto: señala que, en la práctica, la política aprendida de ese
modelo en ese caso es la regla de cambiar de fase en cuanto la máscara
lo permite.
% fuente: trivial: plan2/<esc>_trivial.add.xml (10/3/10/3, desfases en cero)

Todo lo demás es igual que en el plan: las versiones evaluadas, con
sus versiones guardadas y sus entrenamientos elegidos por la validación
en \texttt{timeLoss} (en P6, las del cierre de ese caso), las semillas,
la estadística pareada y la regla de veredicto. El resumen en un número
y el análisis de varianza de la
Sección~\ref{sec:analisis-complementarios} se repiten con el retraso
total y el fijo ajustado desde 10~s. El análisis cubre los seis casos,
P1 a P6.
% fuente: plan2/robustez_veredicto.csv, robustez_sintesis_numero_unico.csv y robustez_sintesis_anova.csv
% (28.09.2026, casos P1/P2/P3/P4/P5/P6); plan2/barrido_malla3_alta.csv (72 filas, avenida 10-35 y transversal
% 10-25; la fila de menor retraso_total_prom, 82.63, es la de menor timeloss_prom, 36.91: 15/10, cero)
% fuente: veredicto.py (--prefijo robustez_; IGUAL = 0.0005, igual_a sobre METRICA y
% espera_prom, marca "(= trivial)"; REFERENCIAS = actuado y trivial); sintesis_plan.py --prefijo robustez_;
% git HEAD:barrer_fijo.py (MALLAS del 25.09: avenida desde 12 s en corredor y 15 s en los demas; en_borde) y
% plan2/cola/logs/barrido_<esc>.log del 25.09 sin AVISO; comentario de MALLAS en barrer_fijo.py (27.09.2026);
% la igualdad de la politica de P3 y el trivial: este anexo, "Que cambia" (antes capitulo4.tex, sec:remedicion)

% Verificacion de la fase 2 (23.10.2026): aqui iba el parrafo "La medicion mas estricta cambia tres cosas" (de
% capitulo5.tex l. 1144-1170 del commit 7181ab8). Su metrica, su barrido y su ciclo minimo estan en los tres
% parrafos de arriba; sus programas, mas abajo, en "Fijo ajustado desde 10 s"; su baja del fijo (21.72 -> 19.91,
% 30.05 -> 28.23 y 20.84 -> 20.78 s/veh, que son de las semillas 1001-1030) va alli una sola vez, junto al
% intervalo y al valor p; lo de P6 tambien, y que las versiones evaluadas son las del plan, en el parrafo anterior.

% Integracion (23.10.2026): el parrafo siguiente es capitulo4.tex l. 1299-1317 del commit 7181ab8 (4.7.5), con su
% fuente (l. 1318-1323). Sin reescribir; salen sus dos ultimas frases ("Ampliar la malla hacia 10 s se decidio
% despues de ver ese resultado, y la especificacion del plan no lo exigia. Por eso el re-barrido forma parte del
% analisis de robustez y no cambia el fijo del resultado principal."), que repiten la entrada. Verificacion de la
% fase 2 (23.10.2026): salen tambien el punto de partida del barrido (ya en "La segunda es el fijo de referencia") y
% el hallazgo de P3 (la decision esta en la entrada y la igualdad, en "Lectura por caso"). Verificacion de las
% tareas 4 y 5 (23.10.2026): la primera frase decia "El barrido del plan no llego a 10 s", que choca con P6, cuyo
% barrido del plan si se amplio a 10 s por la regla de interioridad ("Fijo ajustado desde 10 s", mas abajo; Tabla
% tab:robustez-fijo; conclusiones.tex, conclusion 3); se acota a P1-P5
En P1 a P5, el barrido del plan no llegó a 10~s aunque la regla admitía
el verde mínimo como óptimo (Sección~\ref{sec:protocolo}): ningún
ganador quedó en el borde y la rejilla no se amplió. El retraso en función del
programa tiene al menos dos óptimos locales, el ciclo corto
sincronizado y un ciclo de 46~s con los cruces coordinados en onda
verde, y el ganador del plan estaba en el segundo. La regla solo mira si
el ganador toca el borde de la rejilla: con el ganador en el interior,
el otro óptimo quedaba por debajo del borde inferior y la rejilla nunca
se amplió hacia él.
% fuente: comentario de MALLAS en barrer_fijo.py (27.09.2026); mallas del plan en git (HEAD:barrer_fijo.py:
% avenida 12-30 s en corredor y 15-35 s en los demas); plan2/cola/logs/barrido_<esc>.log del 24.09 sin AVISO;
% igualdad de ippo_vecinos publicada en P3 (base 4042) y el trivial: parciales/eval_red_trivial.csv frente a
% eval_red_ippo_vecinos_s4042.csv, 30 semillas, todas las columnas de trafico iguales (solo difieren
% decisiones y el redondeo de ret_*; revisado 27.09.2026); plan2/ESPEC.md linea 95 (interioridad, 10 s piso
% valido; no pide partir de 10 s)
% Integracion (23.10.2026): el parrafo de entrada del bloque "Tablas del analisis de robustez" del Anexo A
% (anexos.tex l. 136-152 del commit 7181ab8) no se repite aqui: la metrica, el fijo re-barrido desde 10 s, P6 con el
% mismo fijo, el control trivial, las mismas politicas sin reentrenar, la decision tomada despues de ver P3 y que no
% reemplaza al veredicto ya estan en la entrada y en los parrafos de arriba. Su primera frase, la que presenta las
% tablas, va al comienzo de "Tablas por modelo".

\subsection*{Fijo ajustado desde 10~s}

% Integracion (23.10.2026): capitulo4.tex l. 1325-1428 del commit 7181ab8 (4.7.5), con la Tabla tab:robustez-fijo
% (aqui F.1). Sin reescribir; solo la opcion de linea de ordenes de l. 1327 (--metrica retraso_total_prom) se dice en
% palabras. En la pasada del glosario (23.10.2026) el texto y el caption de la tabla pasaron a los nombres nuevos.
% Verificacion de la fase 2 (23.10.2026): "malla" del barrido pasa a "rejilla de valores" y los nombres internos de
% escenario (corredor, corredor_alta, red, red_alta, malla3, malla3_alta) a P1-P6, tambien en la tabla, que pierde la
% columna Escenario (repetia el caso); aqui se dice, una sola vez en el anexo, que P6 no se vuelve a barrer.
\begin{sloppypar}
El barrido desde 10~s usa \texttt{barrer\_fijo.py} con su opción por omisión,
el retraso total como criterio, y la rejilla de los dos verdes desde
10~s en todos los casos: avenida de 10 a 30~s en P1 y de 10 a 35~s en
los demás. Ningún ganador de P1 a P5 quedó en el borde superior de su
rejilla, y ningún barrido tuvo que ampliarla. P6 no se vuelve a barrer.
Su barrido del plan, después de la ampliación a 10~s por la regla de
interioridad (Sección~\ref{sec:referencias-fijas}), cubre la misma
rejilla que el barrido desde 10~s de P5, 72 configuraciones con los dos
verdes desde 10~s, y su ganador por \texttt{timeLoss}, 15~s de avenida
y 10~s de transversal sin desfases, es también el de menor retraso
total. En P6 el fijo del análisis de robustez es, entonces, el mismo
del resultado principal. La Tabla~\ref{tab:robustez-fijo} da el fijo de
los seis casos.
\end{sloppypar}
% fuente: plan2/cola/logs/v21_barrido_<esc>.log (sin AVISO de borde) y plan2/barrido_<esc>.csv de la re-medicion
% (los verdes barridos son exactamente los de barrer_fijo.MALLAS en los cinco casos); P6: plan2/barrido_malla3_alta.csv
% (72 filas, avenida 10-35, transversal 10-25, tres esquemas; minimo de retraso_total_prom 82.63 y de timeloss_prom
% 36.91 en la misma fila); plan2/robustez_veredicto.csv (P6 frente al brazo fijo = eval_malla3_alta_fijo.csv)

\begin{table}[H]
\setstretch{1}
\centering
\scriptsize
\setlength{\tabcolsep}{3pt}
\begin{tabular}{lcclcccc}
\toprule
Caso & \makecell{Verde avenida /\\transversal (s)} & \makecell{Ciclo\\(s)} & Desfases (s) & \makecell{Retraso total\\(s/veh)} & \makecell{\texttt{timeLoss}\\(s/veh)} & Config. & \makecell{Posición\\(av. / tr.)} \\
\midrule
P1 & 20 / 10 & 36 & 0 / 15          & 20.88 (0.22)   & 20.50  & 170 & int. / piso \\
P2 & 25 / 15 & 46 & 0 / 25          & 384.97 (11.58) & 152.13 & 119 & int. / int. \\
P3 & 10 / 10 & 26 & 0 / 25 / 24     & 20.20 (0.19)   & 19.85  & 471 & piso / piso \\
P4 & 10 / 10 & 26 & 0 / 25 / 24     & 26.52 (3.60)   & 26.19  & 471 & piso / piso \\
P5 & 10 / 10 & 26 & todos en 0      & 21.23 (0.10)   & 20.77  & 72  & piso / piso \\
P6 & 15 / 10 & 31 & todos en 0      & 82.63 (0.40)   & 36.91  & 72  & int. / piso \\
\bottomrule
\end{tabular}
\caption[Robustez: fijo ajustado desde 10~s por caso]{Análisis de robustez: fijo ajustado desde
10~s de cada caso, el programa fijo ganador del barrido que parte del verde mínimo (semillas 2001
a 2003), elegido por el retraso total por vehículo
(Ecuación~\ref{eq:retraso-total}). Ciclo: los dos verdes más dos ámbares de 3~s. Desfases: del
primer cruce al último. Retraso total: media y, entre paréntesis, desviación estándar sobre las
tres semillas. \texttt{timeLoss}: media del mismo programa. Config.: configuraciones barridas.
Posición: dónde queda el verde de la avenida y el transversal del ganador en su rejilla; int.:
interior; piso: el verde mínimo de 10~s, que no cuenta como borde. Rejillas: verde de la avenida de
10 a 30~s en P1 y de 10 a 35~s en los demás; verde transversal de 10 a 20~s (10 a 25~s en
P2, P5 y P6); desfase del segundo cruce de 0 a 40~s (0 a 30~s en
P2); en P5 y P6, tres esquemas de desfases. En P6 no hay un barrido
nuevo: el programa es el ganador del barrido del plan, que ya cubre esta rejilla y es también el de menor
retraso total. El fijo del resultado principal es el de la Tabla~\ref{tab:plan2-fijo}.}
\label{tab:robustez-fijo}
\end{table}
% fuente: plan2/barrido_<esc>.csv de la re-medicion del 27.09.2026 (fila de menor retraso_total_prom; columnas
% verde, verde_transversal, offsets, esquema, retraso_total_prom, retraso_total_sd, timeloss_prom; configuraciones
% = filas del csv), tomadas de la version de este capitulo del 27.09 (tab:plan2-fijo v2.1) y de
% scratchpad/escritura/hechos_v21.md (Fijo ganador del re-barrido); ciclo: cabecera de plan2/<esc>.add.xml;
% mallas: barrer_fijo.MALLAS. plan2/barrido_corredor.csv se estaba restaurando el 28.09 y no se releyo; P1 coincide
% con plan2/barrido_corredor_f100.csv (misma configuracion); P6: plan2/barrido_malla3_alta.csv (28.09.2026; fila
% 15/10 cero: retraso_total_prom 82.63, retraso_total_sd 0.4, timeloss_prom 36.91; 72 filas); ciclo 31 s:
% cabecera de plan2/malla3_alta.add.xml

En P1 y P2 el barrido desde 10~s conserva el programa del plan: 20/10~s con
desfase de 15~s en P1 y 25/15~s con desfase de 25~s en P2. En esos dos
casos el ciclo corto queda lejos: el mejor programa con 10~s de verde
en cada fase da 102.71~s/veh de retraso total en P1 y 961.54~s/veh en
P2, casi cinco veces y 2.5 veces el del ganador.

En P3, P4 y P5 el ganador cambia. El barrido del plan había elegido un
ciclo de 46~s (25/15~s en P3 y P5, 30/10~s en P4); con la rejilla que
parte del mínimo gana el ciclo corto de 26~s, con 10~s de verde en cada
fase. En P3 y P4 sus desfases son 0, 25 y 24~s (el tercero es el doble
del segundo, módulo el ciclo), a 1 y 2~s de la sincronía en un ciclo de
26~s. En P5 gana el esquema con todos los desfases en cero, así que el
fijo ajustado desde 10~s de P5 es el mismo programa que el ciclo
mínimo. El esquema en ajedrez, que ganaba el barrido del plan, queda
segundo, con 25/15~s y 21.28~s/veh.

La ventaja del nuevo ganador sobre el del plan, medida con el retraso
total en las mismas semillas de ajuste, no es igual de clara en los
tres casos. En P3 es de 1.67~s/veh, casi nueve veces la desviación
estándar del ganador entre esas semillas (0.19). En P4 es de
1.84~s/veh, menos que esa desviación (3.60), y en P5, de 0.05~s/veh,
frente a 0.10. En estos dos casos los dos óptimos quedan dentro de la
variación entre las tres semillas de ajuste, y el barrido elige el de
menor media, como fija la regla.
% fuente: plan2/barrido_<esc>.csv de la re-medicion (retraso_total_prom y retraso_total_sd): corredor 10/10
% mejor 0/25 = 102.71 frente a 20.88; corredor_alta 10/10 mejor 0/10 = 961.54 frente a 384.97; ganador del
% barrido del plan (git, HEAD:plan2/<esc>.add.xml) en la malla nueva: red 25/15 0/23/0 = 21.87 (sd 0.28) frente
% a 20.20 (sd 0.19); red_alta 30/10 0/20/0 = 28.36 (sd 3.04) frente a 26.52 (sd 3.60); malla3 25/15 ajedrez
% 0/23/... = 21.28 (sd 0.31) frente a 21.23 (sd 0.10); los tres son el mejor programa con algun verde de mas
% de 10 s en su csv. Igualdad del fijo de P5 y el trivial: plan2/malla3.add.xml y malla3_trivial.add.xml
% (10/3/10/3, offset 0 en los nueve)

La Tabla~\ref{tab:robustez-fijo} da el retraso de cada programa en las
semillas de ajuste, 2001 a 2003. En las de evaluación, 1001 a 1030, y
con el \texttt{timeLoss}, la métrica del resultado principal, la
diferencia es clara en P3 y también en P4. Frente al fijo del plan, el
fijo ajustado desde 10~s baja el \texttt{timeLoss} de 21.72 a
19.91~s/veh en P3, un 8.35\,\% (intervalo de $-8.87$ a $-7.80$\,\%), y
de 30.05 a 28.23~s/veh en P4, un 6.06\,\% ($-7.54$ a $-4.61$\,\%), con
$p < 0.001$ en los dos. En P5 el cambio, de 20.84 a 20.78~s/veh
($-0.28$\,\%, de $-0.58$ a $0.03$\,\%, $p = 0.090$), no es
significativo. En P3 esa diferencia es del mismo tamaño que la mejora
de IPPO con mensajes frente al fijo del plan ($-8.8$\,\%), y en P4 la
supera ($-3.6$\,\%).
% fuente: scratchpad/escritura/hechos_2509.md (Fijo del 25.09 frente al fijo re-barrido, timeLoss, 1001-1030:
% P3 -8.35 % [-8.87, -7.80] p=1.73e-06; P4 -6.06 % [-7.54, -4.61] p=8.01e-08; P5 -0.28 % [-0.58, 0.03]
% p=0.0896); medias en 1001-1030 (verificadas el 23.10.2026): brazo fijo frente a fijo_2509, metrica timeloss_prom,
% columnas media, delta_pct, delta_ic95 y p_wilcoxon de plan2/evaluacion_<esc>_resumen_fijo_2509.csv (red 19.9053
% frente a 21.7187; red_alta 28.2317 frente a 30.0517; malla3 20.7803 frente a 20.8387);
% plan2/veredicto_congelado_0355.csv (ippo_vecinos delta_fijo P3 -8.8 %, P4 -3.6 %)

\subsection*{Ejecución}

% Integracion (23.10.2026): capitulo4.tex l. 1430-1460 y 1484-1529 del commit 7181ab8 (4.7.5); las l. 1462-1482 van a
% "Comprobaciones". Sin reescribir; solo cambian: las opciones de linea de ordenes de l. 1500-1505 (--prefijo
% robustez_, --casos y --metrica timeloss_prom sin --fijo) se dicen en palabras, y el parrafo del antivirus va en
% sloppypar, porque en el Cap. IV se salia 27 pt del margen por los identificadores largos.
% Verificacion de la fase 2 (23.10.2026; plan del 14.10, seccion 3.1: la bitacora del proceso se quita): el apartado
% se condenso. Salen la funcion agregar_v21 y sus prioridades, el detalle de las siete evaluaciones de IPPO con
% penalizacion que corrieron dos veces, el antivirus y las opciones de los programas que escriben las tablas; queda
% lo que se corrio, que dos barridos se repitieron sin cambiar ningun numero y el costo. Las fuentes siguen abajo.
En P1 a P5 la medición más estricta tiene tres grupos de trabajos. El
primero es el barrido desde 10~s del fijo de cada caso, que instala el
ganador como programa fijo del caso. El segundo es la evaluación del
fijo, del actuado y del ciclo mínimo con las semillas 1001 a 1030. El
tercero es la evaluación de todas las políticas ya entrenadas, los tres
entrenamientos de cada modelo y caso, con las mismas semillas, para
obtener el retraso total, la espera para entrar y los no insertados,
que el plan no registraba. Con el código nuevo corrieron también las
siete evaluaciones de IPPO con penalización de \textit{spillback} que
habían quedado pendientes (Sección~\ref{sec:spill-vecinos}). La
elección de la versión guardada de cada entrenamiento y de la versión
evaluada no se rehízo: sigue la validación en \texttt{timeLoss} del
plan (Sección~\ref{sec:protocolo}). Dos barridos se interrumpieron por
un error al escribir un archivo de salida y se repitieron con la misma
orden, sin que cambiara ningún número de la simulación. Las tablas las
escriben los mismos programas del resultado principal, con el retraso
total y el fijo ajustado desde 10~s.
% fuente: plan2/cola/cola.out (eval_<esc>_ippo_spill_s<b> de red, red_alta y malla3 entre las 18:35 y las
% 19:03; v21_eval_<esc>_ippo_spill_s<b> despues, por ejemplo v21_eval_red_ippo_spill_s4042 de 20:03 a 20:06)
% y plan2/cola/estado.json (tiempos de los dos identificadores); barridos interrumpidos: FALLO de v21_barrido_red y
% v21_barrido_red_alta a las 19:00 y 19:15 en plan2/cola/cola.out, repetidos con exito; docstrings de
% barrer_fijo.correr_sumo (5 intentos, 3 s; "tumbo dos barridos"), comun.borrar_tripinfo (5 intentos, 1 s) y
% comun.iniciar_sumo (24.09.2026); tablas: docstrings de veredicto.py, sintesis_plan.py y plan2/tablas_cap5.py
% (--prefijo robustez_; 28.09.2026); plan2/robustez_veredicto.csv (28.09.2026, P1 a P6): sus filas de P1 a P5
% coinciden con plan2/veredicto_v21_respaldo.csv, donde P6 aun figuraba sin evaluacion

La medición más estricta no entrena, así que su costo lo fijan los
barridos y las evaluaciones. En P1 a P5 son 137 trabajos: los cinco
barridos del fijo y 132 evaluaciones, las del fijo, el actuado y el
ciclo mínimo de cada caso y las de 117 políticas, tres entrenamientos
por modelo y caso. Terminaron todos, sin ninguno fallido al cierre. Sus
duraciones suman 16.4 horas: 2.0 en los barridos y 14.4 en las
evaluaciones; como cada barrido corre tres simulaciones a la vez, el
tiempo de CPU puede llegar a 20.4 horas. P6 no suma trabajos a la
medición más estricta: sus evaluaciones, incluida la del ciclo mínimo,
corrieron con el código que ya registra el retraso total
(Sección~\ref{sec:p6-demanda}).
% fuente: plan2/cola/estado.json (tiempos, en s, de los trabajos v21_*; hechos; fallidos vacio):
% v21_ 137 trabajos, 59 086 s (barridos 5, 7173 s; evaluaciones 132, 51 913 s); cota de CPU = evaluaciones + 3 x
% barridos (--procesos 3 en plan2/cola/trabajos.jsonl); plan2/cola/cola.out (INICIO v21_barrido_corredor
% 18:35:22; ultimo OK v21_ 20:55:26; FALLO de v21_barrido_red y v21_barrido_red_alta a las 19:00 y 19:15,
% repetidos con exito: son los dos barridos que interrumpio el antivirus); las siete
% eval_<esc>_ippo_spill_s<b> pendientes del plan corrieron con su identificador original entre las 18:35 y las
% 19:03 (2688 s en total) y despues otra vez como v21_eval_*; no se suman a los 137

\subsection*{Veredictos con la medición más estricta}

La Tabla~\ref{tab:robustez-veredicto-matriz} da el veredicto con estos
cambios, y la Tabla~\ref{tab:robustez-plan2-resumen}, el retraso total
y el cambio frente al fijo de cada modelo. La nota ``igual al ciclo mínimo''
señala las versiones evaluadas que dan el mismo retraso total y la
misma espera detenida que el ciclo mínimo.

% fuente: plan2/robustez_veredicto.csv (= plan2/veredicto_v21_respaldo.csv, la re-medicion del 27.09.2026); se
% regenera con "python veredicto.py --prefijo robustez_" desde demo_rl/ (con los seis casos; hasta el 28.09 se
% corria con --casos P1-P5)
% Generado por veredicto.py a partir de plan2/. No editar a mano.
\begin{table}[H]
\setstretch{1}
\centering
\scriptsize
\setlength{\tabcolsep}{3pt}
\begin{tabular}{lcccccc}
\toprule
Modelo & \makecell{P1\\Corredor\\nominal} & \makecell{P2\\Corredor\\alta} & \makecell{P3\\Red con\\rotonda\\nominal} & \makecell{P4\\Red con\\rotonda\\alta} & \makecell{P5\\Malla $3 \times 3$\\nominal} & \makecell{P6\\Malla $3 \times 3$\\alta} \\
\midrule
Actuado & EMPATA & FALLA & FALLA & FALLA & FALLA & FALLA \\
Trivial & FALLA & FALLA & MEJORA & EMPATA & EMPATA & FALLA \\
Q-learning & FALLA & FALLA & FALLA & PENDIENTE & FALLA & FALLA \\
SARSA & FALLA & FALLA & FALLA & FALLA & FALLA & PENDIENTE \\
Q-learning (colas) & FALLA & PENDIENTE & FALLA & FALLA & PENDIENTE & PENDIENTE \\
DQN & MEJORA & FALLA & FALLA & FALLA & PENDIENTE & PENDIENTE \\
IPPO local & EMPATA & MEJORA & FALLA & FALLA & FALLA & FALLA \\
IPPO vecinos & FUNCIONA & FALLA & MEJORA* & FALLA & FALLA & PENDIENTE \\
IPPO spill & FUNCIONA & FALLA & = vecinos & = vecinos & = vecinos & PENDIENTE \\
MAPPO & FUNCIONA & FALLA & PENDIENTE & PENDIENTE & EMPATA* & PENDIENTE \\
\bottomrule
\end{tabular}
\caption[Robustez: veredicto por caso y modelo]{Robustez. Veredicto de cada modelo en cada caso, con la regla de la Tabla~\ref{tab:veredictos}: retraso total (timeLoss más la espera para entrar a la red) de la versión evaluada frente al fijo re-barrido y al actuado, semillas 1001 a 1030. IPPO local, vecinos y spill: IPPO sin mensajes, con mensajes y con penalización de \textit{spillback}; DQN: Double DQN. = vecinos: la penalización no se activa e IPPO spill lleva el veredicto de IPPO vecinos. PENDIENTE: faltan datos o el entrenamiento de la versión evaluada no se estabilizó. El actuado y el trivial (verdes mínimos, desfases en 0) se juzgan frente al fijo. *: la política da el mismo retraso y la misma espera que el trivial.}
\label{tab:robustez-veredicto-matriz}
\end{table}


% fuente: plan2/robustez_veredicto.csv y fila fijo (retraso_total_prom) de plan2/evaluacion_<esc>_resumen.csv;
% tabla escrita por "python tablas_cap5.py --prefijo robustez_" desde plan2/ (con los seis casos). Verificacion de
% las tareas 4 y 5 (23.10.2026): flota ([!tbp]); en [H], detras de la matriz, que flota, las dos quedaban en la misma
% pagina y el anexo se salia 261 pt por abajo. El retraso de cada modelo se redondea una sola vez, desde la media de
% evaluacion_<esc>_resumen.csv, como en las tablas por modelo (antes DQN en P1 daba 20.93 aqui y 20.94 en la Tabla
% tab:robustez-veredicto-bases-dqn)
% Generado por plan2/tablas_cap5.py. No editar a mano.
\begin{table}[H]
\setstretch{1}
\centering
\scriptsize
\setlength{\tabcolsep}{2.5pt}
\begin{tabular}{lcccccc}
\toprule
Modelo & \makecell{P1\\Corredor\\nominal} & \makecell{P2\\Corredor\\alta} & \makecell{P3\\Red con\\rotonda\\nominal} & \makecell{P4\\Red con\\rotonda\\alta} & \makecell{P5\\Malla $3 \times 3$\\nominal} & \makecell{P6\\Malla $3 \times 3$\\alta} \\
\midrule
\makecell[l]{Fijo ajustado\\desde 10~s\\(s/veh)} & \makecell{21.67\\$\pm$ 0.26} & \makecell{394.27\\$\pm$ 3.67} & \makecell{20.26\\$\pm$ 0.11} & \makecell{28.57\\$\pm$ 1.11} & \makecell{21.25\\$\pm$ 0.05} & \makecell{82.19\\$\pm$ 0.16} \\
\midrule
Actuado & \makecell{21.45\\$-$1.0 EMP} & \makecell{445.46\\+13.0 FAL} & \makecell{21.84\\+7.8 FAL} & \makecell{30.37\\+6.3 FAL} & \makecell{25.84\\+21.6 FAL} & \makecell{109.18\\+32.8 FAL} \\
Ciclo mínimo & \makecell{106.26\\+390.3 FAL} & \makecell{1055.41\\+167.7 FAL} & \makecell{20.17\\$-$0.5 MEJ} & \makecell{28.29\\$-$1.0 EMP} & \makecell{21.25\\0.0 EMP} & \makecell{276.72\\+236.7 FAL} \\
Q-learning & \makecell{23.64\\+9.1 FAL} & \makecell{498.52\\+26.4 FAL} & \makecell{24.76\\+22.2 FAL} & \makecell{33.36\\+16.8\\PEN (FAL)} & \makecell{30.16\\+41.9 FAL} & \makecell{92.84\\+13.0 FAL} \\
SARSA & \makecell{24.64\\+13.7 FAL} & \makecell{474.14\\+20.3 FAL} & \makecell{24.50\\+21.0 FAL} & \makecell{33.52\\+17.4 FAL} & \makecell{29.65\\+39.5 FAL} & \makecell{90.95\\+10.7\\PEN (FAL)} \\
\makecell[l]{Q-learning\\(colas)} & \makecell{23.93\\+10.4 FAL} & \makecell{559.95\\+42.0\\PEN (FAL)} & \makecell{24.16\\+19.3 FAL} & \makecell{33.34\\+16.7 FAL} & sin datos & sin datos \\
DQN & \makecell{20.93\\$-$3.4 MEJ} & \makecell{461.12\\+17.0 FAL} & \makecell{21.68\\+7.0 FAL} & \makecell{29.37\\+2.8 FAL} & \makecell{23.62\\+11.2\\PEN (FAL)} & \makecell{86.33\\+5.0\\PEN (FAL)} \\
\makecell[l]{IPPO sin\\mensajes} & \makecell{21.56\\$-$0.5 EMP} & \makecell{387.46\\$-$1.7 MEJ} & \makecell{22.57\\+11.4 FAL} & \makecell{31.25\\+9.4 FAL} & \makecell{26.70\\+25.6 FAL} & \makecell{84.68\\+3.0 FAL} \\
\makecell[l]{IPPO con\\mensajes} & \makecell{20.22\\$-$6.7 FUN} & \makecell{400.68\\+1.6 FAL} & \makecell{20.17\\$-$0.5 MEJ\\igual al\\ciclo mínimo} & \makecell{29.30\\+2.6 FAL} & \makecell{21.68\\+2.0 FAL} & \makecell{83.05\\+1.1\\PEN (EMP)} \\
\makecell[l]{IPPO con\\penalización} & \makecell{20.57\\$-$5.1 FUN} & \makecell{421.45\\+6.9 FAL} & \makecell{igual a IPPO\\con mensajes} & \makecell{igual a IPPO\\con mensajes} & \makecell{igual a IPPO\\con mensajes} & \makecell{83.58\\+1.7\\PEN (EMP)} \\
MAPPO & \makecell{20.16\\$-$7.0 FUN} & \makecell{402.56\\+2.1 FAL} & \makecell{20.37\\+0.5\\PEN (EMP)} & \makecell{29.78\\+4.3\\PEN (FAL)} & \makecell{21.25\\0.0 EMP\\igual al\\ciclo mínimo} & \makecell{86.82\\+5.6\\PEN (FAL)} \\
\bottomrule
\end{tabular}
\caption[Robustez: retraso y cambio frente al fijo por caso y modelo]{Robustez. Resumen del plan de pruebas por caso y modelo, sobre las semillas de evaluación 1001 a 1030, con el retraso total (Ecuación~\ref{eq:retraso-total}). Primera fila: el del fijo ajustado desde 10~s (s/veh), media $\pm$ semiamplitud del IC95\,\%. Demás filas: el del modelo (s/veh), es decir, de su versión evaluada, y debajo $\Delta_f$ (\%) con el veredicto abreviado: FUN, funciona; MEJ, mejora; EMP, empata; FAL, falla; PEN, pendiente porque el entrenamiento de la versión evaluada no converge, con el veredicto provisional entre paréntesis. Igual al ciclo mínimo: la política da el mismo retraso total y la misma espera que el ciclo mínimo. Igual a IPPO con mensajes: la penalización de \textit{spillback} no se activa y la política es la de IPPO con mensajes. Sin datos: Q-learning (colas) no se corrió en P5 ni en P6. Los intervalos, los valores $p$ y $\Delta_a$ están en las tablas por modelo.}
\label{tab:robustez-plan2-resumen}
\end{table}


Con la medición estricta, ningún aprendiz funciona fuera de P1. En P1
funcionan los mismos tres modelos que en el veredicto pre-registrado.
IPPO con mensajes mejora al fijo en P3 solo porque su política es el
ciclo mínimo, falla en P2, P4 y P5, y en P6 queda pendiente, con empate
provisional. MAPPO falla en P2 y empata en P5 con una política que
también es el ciclo mínimo, y en P3, P4 y P6 queda pendiente. IPPO sin
mensajes es el único aprendiz que mejora al fijo en P2 ($-1.7$\,\%). El
actuado falla en P2 a P6. De los 45 veredictos de modelos y actuado de
P1 a P5, contando los provisionales, 26 no cambian. Diez cambian por el
fijo ajustado desde 10~s, todos en P3 y P4; siete por la métrica, todos
en P2; y dos, los de IPPO con mensajes e IPPO con penalización de
\textit{spillback} en P5, por un borde estadístico. Para separar las
dos causas, la regla se aplicó también con el \texttt{timeLoss} frente
al fijo ajustado desde 10~s. En P6 no cambia ninguno de los nueve
veredictos. La Figura~\ref{fig:robustez-mejora-caso} compara, para
IPPO con mensajes y MAPPO, el cambio frente al fijo del veredicto
pre-registrado con el de la medición más estricta.
% fuente: plan2/robustez_veredicto.csv; atribucion de causas: tabla_antes_despues.md de la auditoria del
% 27.09.2026 (25.09 -> timeLoss con el fijo re-barrido [fijo] -> retraso total con el fijo re-barrido [metrica];
% conteo: sin cambio 26, fijo 10, metrica 7, borde estadistico 2, trivial 5 celdas nuevas); P6: columna veredicto
% de malla3_alta en plan2/veredicto.csv y plan2/robustez_veredicto.csv, igual en las nueve filas de modelos y actuado

% Verificacion de la fase 2 (23.10.2026): flota ([!tbp]); en [H] no cabia al pie de la pagina y dejaba en blanco
% mas de media pagina
\begin{figure}[!tbp]
\centering
\includegraphics[width=\textwidth]{fig_robustez_mejora_caso.png}
\caption[Robustez: cambio frente al fijo, veredicto y medición más estricta]{Robustez. Cambio del retraso frente al fijo (\%) de las versiones evaluadas de IPPO con mensajes (azul) y MAPPO (naranja) en P1 a P6; P6 solo tiene el valor de la medición más estricta, porque el veredicto congelado no incluía ese caso (su cambio con el \texttt{timeLoss} está en la Tabla~\ref{tab:plan2-resumen}). Círculo hueco: veredicto pre-registrado, con el \texttt{timeLoss} y el fijo ajustado (los valores de la Tabla~\ref{tab:plan2-resumen}); a diferencia de la Figura~\ref{fig:plan2-mejora-caso}, aquí el hueco no indica un veredicto PENDIENTE. Punto relleno: medición más estricta, con el retraso total y el fijo ajustado desde 10~s; la barra vertical es su IC95\,\% por \textit{bootstrap} pareado sobre las semillas 1001 a 1030. La línea tenue une los dos valores de una misma política. Línea roja discontinua: meta de la tesis, $-25$\,\%. Bajo cada caso, su número de semáforos, los carriles por sentido de la avenida y la demanda total (veh/h). Un valor negativo indica menos retraso que el fijo.}
\label{fig:robustez-mejora-caso}
\end{figure}
% fuente: figura escrita por "python sintesis_plan.py --prefijo robustez_" desde
% plan2/robustez_veredicto.csv (delta_fijo, ic_fijo) y plan2/veredicto_congelado_0355.csv (delta_fijo); regenerada
% con P6 el 28.09.2026: P6 solo tiene el punto de la re-medicion, porque el congelado del 25.09 no lo incluye.
% Verificacion de la fase 2 (23.10.2026): el caption aclara que aqui el hueco no es PENDIENTE (en la Figura
% fig:plan2-mejora-caso si lo es) y dice "cambio del retraso frente al fijo" en lugar de Delta_f

\subsection*{Lectura por caso}

\noindent\textbf{P1: sin cambio.} El programa fijo es el mismo, y la
espera para entrar es de 0.37~s/veh con el fijo, el actuado y las
versiones evaluadas de IPPO con mensajes, Double DQN y MAPPO, así que
el retraso total solo suma una constante. IPPO con mensajes queda en $-6.7$\,\% (frente a $-6.8$\,\%
en el veredicto pre-registrado), MAPPO en $-7.0$\,\% e IPPO con
penalización de \textit{spillback} en $-5.1$\,\%, y ningún veredicto
cambia. El ciclo mínimo falla por mucho ($+390.3$\,\%): en el
corredor, con la demanda nominal, el ciclo corto no alcanza.
% fuente: plan2/robustez_veredicto.csv (P1); plan2/evaluacion_corredor_resumen.csv (retraso_insercion_prom
% 0.37 en fijo, actuado, ippo_vecinos_s4042, mappo_s42 y dqn_s4042). 23.10.2026 (entrega 3, pedido 2b): antes
% decia "con un carril por sentido y la demanda nominal"; el plan no aisla el efecto de los carriles

\noindent\textbf{P2: la métrica.} El programa fijo tampoco cambia;
cambia lo que se mide. Con el fijo, cada vehículo espera en promedio
más para entrar a la red (239.6~s) que dentro de ella (154.7~s), y el
retraso total es de 394.3~s/veh. La política de IPPO con mensajes sigue con
un \texttt{timeLoss} un $23.3$\,\% menor que el del fijo, pero su
espera para entrar sube a 282.0~s/veh y su retraso total queda un
$1.6$\,\% por encima, una falla. El actuado baja el \texttt{timeLoss}
un $19.8$\,\% y sube el retraso total un $13.0$\,\%, con 321.3~s/veh
de espera para entrar. Estos controles despachan mejor a los vehículos
que ya están en la red y dejan crecer la cola en las entradas, donde
ni el \texttt{timeLoss} ni la recompensa la ven
(Sección~\ref{sec:formulacion}). Ningún vehículo queda sin entrar ni
sin terminar con el fijo, el actuado y esas políticas, así que su
retraso total es el valor medido y no una cota. IPPO sin mensajes es
el único aprendiz que baja el retraso total ($-1.7$\,\%, $p = 0.001$).
% fuente: plan2/evaluacion_corredor_alta_resumen.csv (fijo: timeloss_prom 154.7013, retraso_insercion_prom
% 239.5657, retraso_total_prom 394.2653; retraso_insercion_prom de ippo_vecinos_s2042 282.0237 y actuado 321.32;
% no_insertados y no_terminados 0.0); plan2/robustez_veredicto.csv (P2)

\noindent\textbf{P3 y P4: el fijo.} En estos casos la métrica casi no
importa, porque la espera para entrar es en promedio de 0.35~s/veh o
menos con todos los controles. Lo que cambia es la referencia. El
ciclo mínimo, que la rejilla del primer barrido no incluía, tiene un
retraso total un $8.6$\,\% menor que el del fijo del primer barrido en
P3 y un $6.9$\,\% menor en P4. La versión evaluada de IPPO con mensajes
en P3 da, en las treinta semillas y en todas las métricas de tráfico,
lo mismo que el ciclo mínimo: su $-8.8$\,\% del veredicto pre-registrado era la
ventaja del ciclo corto sobre el largo, no la de una coordinación
aprendida. Frente al fijo ajustado desde 10~s esa política mejora
apenas ($-0.5$\,\%, $p = 0.013$) y lleva la nota ``igual al ciclo
mínimo''. En P4, donde el ciclo mínimo empata con el fijo ajustado
desde 10~s ($-1.0$\,\%, $p = 0.22$), IPPO con mensajes, Double DQN y
MAPPO quedan por encima del fijo, entre $+2.6$ y $+4.3$\,\%.
% fuente: plan2/evaluacion_red_resumen.csv y plan2/evaluacion_red_alta_resumen.csv (retraso_insercion_prom 0.35 y
% 0.34); trivial frente al fijo del 25.09 en retraso total: contra_trivial.txt de la auditoria del 27.09.2026
% (P3 -8.6 % [-9.2, -8.1], P4 -6.9 % [-8.1, -5.8]); igualdad con el trivial: igual_trivial.txt de la auditoria
% (red ippo_vecinos_s4042 = trivial en 30/30 semillas) y parciales/eval_red_trivial.csv frente a
% eval_red_ippo_vecinos_s4042.csv (todas las columnas de trafico iguales; revisado 27.09.2026; antes en el parrafo
% de los dos optimos de "Que cambia"); plan2/robustez_veredicto.csv (P3 y P4)

\noindent\textbf{P5: un borde estadístico.} El fijo ajustado desde
10~s es el ciclo mínimo, y su \texttt{timeLoss} apenas difiere del fijo
del primer barrido. IPPO con mensajes pasa de empate a falla: su
cambio frente al fijo es de $+2.0$\,\%, con un intervalo de $+0.7$ a
$+3.6$\,\% y $p = 0.028$. Su política da el mismo resultado que el ciclo mínimo en
24 de las 30 semillas; en las otras seis
su \texttt{timeLoss} queda entre 1.0 y 2.8~s/veh por encima. La prueba
de Wilcoxon descarta las diferencias nulas y decide con esas seis,
todas positivas. Con la variante que reparte los ceros entre los dos
signos, el valor $p$ sería 0.070 y el veredicto, empate. La versión
evaluada de MAPPO es el ciclo mínimo y empata. El veredicto de IPPO con
mensajes en P5 es, así, frágil con cualquiera de los dos fijos.
% fuente: plan2/robustez_veredicto.csv (P5); p5_semillas.txt de la auditoria del 27.09.2026 (diferencias por
% semilla de ippo_vecinos_s2042 frente al fijo re-barrido: 24 nulas, seis entre +0.97 y +2.78 s/veh);
% p5_wilcoxon_ceros.txt (wilcox 0.0277, zsplit 0.0702); igual_trivial.txt (malla3 ippo_vecinos_s2042 = trivial en
% 24/30, mappo_s2042 en 30/30)

\noindent\textbf{P6: el mismo fijo y la misma espera para entrar.} El
fijo es el mismo con las dos métricas, y la espera para entrar a la
red es de 45.46 a 45.65~s/veh con el fijo, el actuado y todas las
versiones evaluadas, así que el retraso total suma a cada uno casi la
misma constante; solo el ciclo mínimo llega a unos 99~s/veh. Con el fijo, el retraso
total es de 82.19~s/veh, de los que 45.46 son de espera para entrar.
Los veredictos son los mismos que con el \texttt{timeLoss}, y los
cambios se reducen a cerca de la mitad porque el denominador crece.
IPPO con mensajes queda en $+1.1$\,\% (de $-0.6$ a $+2.8$\,\%,
$p = 0.596$) e IPPO con penalización de \textit{spillback} en
$+1.7$\,\%, los dos con un empate provisional. MAPPO ($+5.6$\,\%), Double DQN
($+5.0$\,\%) y SARSA ($+10.7$\,\%) fallan de forma provisional, e IPPO
sin mensajes ($+3.0$\,\%), Q-learning ($+13.0$\,\%) y el actuado
($+32.8$\,\%) fallan. El ciclo mínimo falla por mucho
($+236.7$\,\%): con la demanda doble, un ciclo de 26~s no alcanza,
como tampoco alcanzaba en los corredores.
% fuente: plan2/robustez_veredicto.csv (P6: delta_fijo, ic_fijo, p_fijo y veredicto); retraso_total_prom y
% retraso_insercion_prom de fijo, actuado, ippo_vecinos_s42 y demas bases publicadas en
% plan2/evaluacion_malla3_alta_resumen.csv (fijo 82.1873 y 45.4617; actuado y
% publicadas entre 45.50 y 45.65, maximo ql_s4042 45.6483; el trivial 99.185, columna media de retraso_insercion_prom,
% verificado el 23.10.2026; antes solo en el bloque "Tablas del analisis de robustez" del Anexo A, anexos.tex l. 171-172
% del commit 7181ab8: "llega a 99 s/veh"); veredictos de P6 con timeLoss: plan2/veredicto.csv

\noindent\textbf{Frente al actuado.} Con la medición estricta, IPPO
con mensajes sigue superando al actuado en los seis casos: en un
$5.7$\,\% en P1, un $10.1$\,\% en P2, un $7.6$\,\% en P3, un
$3.5$\,\% en P4, un $16.1$\,\% en P5 y un $23.9$\,\% en P6. En P2 la
ventaja crece respecto del resultado principal, porque el actuado deja más cola fuera
de la red que la política.
% fuente: plan2/robustez_veredicto.csv (delta_act de ippo_vecinos: -5.7, -10.1, -7.6, -3.5, -16.1, -23.9)

\subsection*{Resumen en un número y reparto de la variación}

La Tabla~\ref{tab:robustez-numero-unico} repite el resumen en un
número. Sobre los seis casos, IPPO con mensajes queda en $-0.03$\,\%
frente al fijo ajustado desde 10~s, con un intervalo de $-0.56$ a
$+0.53$\,\% que incluye el cero, y en $-11.45$\,\% frente al actuado
(de $-11.95$ a $-10.95$\,\%). MAPPO queda en $+0.85$ y $-10.67$\,\%.
Sobre P1 a P5, IPPO con mensajes quedaba en $-0.24$ y $-8.71$\,\%, y
MAPPO en $-0.09$ y $-8.57$\,\%. Frente al resultado principal ($-6.83$
y $-14.35$\,\% con los seis casos), la ventaja frente al fijo
desaparece y la ventaja frente al actuado se mantiene. Sin P1, IPPO con
mensajes pasa a
$+1.36$\,\% frente al fijo. El reparto de la variación también cambia:
con todos los aprendices, el modelo explica el $55.5$\,\%, el caso el
$23.5$\,\%, la interacción el $14.9$\,\% y el entrenamiento el
$6.1$\,\%, al revés que en la Tabla~\ref{tab:plan2-anova}, donde el
caso explica el $53.7$\,\% y el modelo el $36.5$\,\%. Con los cuatro
de redes neuronales el caso sigue siendo el primer factor, pero baja
del $79.5$ al $36.6$\,\%, y la interacción sube al $30.4$\,\%. Sobre
P1 a P5, con todos los aprendices, el modelo explicaba el $59.4$\,\% y
el caso el $23.1$\,\%.
% fuente: plan2/robustez_sintesis_numero_unico.csv (P1-P6: ippo_vecinos -0.03 [-0.56, 0.53], -11.45 [-11.95,
% -10.95], sin_P1 1.36; mappo 0.85 y -10.67; P1-P5, columna sin_P6: ippo_vecinos -0.24 y -8.71, mappo -0.09 y
% -8.57); plan2/robustez_sintesis_anova.csv (P1-P6: Todos 23.5/55.5/14.9/6.1; Profundos 36.6/20.5/30.4/12.5);
% P1-P5 (caso 23.1, modelo 59.4, interaccion 11.4, base 6.1): plan2/robustez_p1p5_sintesis_anova.csv y
% robustez_p1p5_sintesis_numero_unico.csv ("python sintesis_plan.py --casos P1-P5 --prefijo robustez_p1p5_ --tesis
% <carpeta aparte>", 28.09.2026); principal:
% plan2/sintesis_anova.csv y plan2/sintesis_numero_unico.csv

% fuente: python sintesis_plan.py --prefijo robustez_ (desde demo_rl/, con los seis casos)
% -> plan2/robustez_sintesis_numero_unico.csv. Verificacion de la fase 2 (23.10.2026): la tabla va despues del
% parrafo que la presenta; antes iba bajo el titulo y, en [H], no cabia en la pagina y dejaba en blanco casi la mitad
% Generado por sintesis_plan.py. No editar a mano.
\begin{table}[H]
\setstretch{1}
\centering
\scriptsize
\setlength{\tabcolsep}{3pt}
\begin{tabular}{lccrrrrrr}
\toprule
 & \makecell{Frente al fijo\\(\%) [IC95]} & \makecell{Frente al actuado\\(\%) [IC95]} & \multicolumn{6}{c}{Frente al fijo, sin el caso} \\
\cmidrule(lr){4-9}
Modelo & & & P1 & P2 & P3 & P4 & P5 & P6 \\
\midrule
IPPO con mensajes & \makecell[c]{$-$0.03\\{}[$-$0.56, +0.53]} & \makecell[c]{$-$11.45\\{}[$-$11.95, $-$10.95]} & +1.36 & $-$0.36 & +0.06 & $-$0.54 & $-$0.44 & $-$0.24 \\
MAPPO & \makecell[c]{+0.85\\{}[+0.25, +1.44]} & \makecell[c]{$-$10.67\\{}[$-$11.29, $-$10.05]} & +2.49 & +0.60 & +0.91 & +0.18 & +1.02 & $-$0.09 \\
DQN & \makecell[c]{+6.40\\{}[+5.88, +6.93]} & \makecell[c]{$-$5.75\\{}[$-$6.31, $-$5.19]} & +8.48 & +4.41 & +6.28 & +7.13 & +5.48 & +6.68 \\
IPPO sin mensajes & \makecell[c]{+7.49\\{}[+7.07, +7.91]} & \makecell[c]{$-$4.79\\{}[$-$5.30, $-$4.26]} & +9.16 & +9.43 & +6.72 & +7.11 & +4.18 & +8.40 \\
Q-learning & \makecell[c]{+21.12\\{}[+20.30, +21.98]} & \makecell[c]{+7.28\\{}[+6.44, +8.17]} & +23.68 & +20.08 & +20.90 & +22.00 & +17.33 & +22.82 \\
SARSA & \makecell[c]{+20.08\\{}[+19.42, +20.76]} & \makecell[c]{+6.37\\{}[+5.74, +6.99]} & +21.39 & +20.04 & +19.90 & +20.63 & +16.52 & +22.06 \\
Actuado & \makecell[c]{+12.89\\{}[+12.23, +13.54]} & -- & +15.90 & +12.87 & +13.94 & +14.26 & +11.22 & +9.28 \\
Ciclo mínimo & \makecell[c]{+87.58\\{}[+82.30, +92.43]} & -- & +54.78 & +74.70 & +112.92 & +113.14 & +112.72 & +66.86 \\
\bottomrule
\end{tabular}
\caption[Robustez: resumen en un número]{Robustez. Resumen en un número: media geométrica, sobre los casos, del cociente entre el retraso total del modelo (versión evaluada) y el de la referencia, expresada como cambio \%; negativo es mejor. IC95\,\% bootstrap pareado por semilla dentro de cada caso (semillas 1001 a 1030). Las últimas columnas repiten el índice frente al fijo dejando fuera un caso.}
\label{tab:robustez-numero-unico}
\end{table}


La Tabla~\ref{tab:robustez-veredicto-trivial}, al final de este anexo,
da el veredicto del ciclo mínimo
frente al fijo ajustado desde 10~s. Falla en los dos corredores ($+390.3$ y
$+167.7$\,\%), mejora al fijo en P3 ($-0.5$\,\%), empata en P4 y P5,
y vuelve a fallar por mucho en P6 ($+236.7$\,\%). En P5 es el mismo
programa que el fijo.
% fuente: plan2/robustez_veredicto.csv (fila trivial); la tabla robustez_veredicto_trivial.tex se incluye una sola
% vez, en "Tablas por modelo", al final de este anexo, para no duplicar tab:robustez-veredicto-trivial
% Entrega 3, fase 2: antes decia "en el Anexo A con las demas tablas por modelo"; desde la integracion del 23.10.2026
% el bloque de tablas por modelo esta al final de este anexo
% Reduccion del Anexo F (23.10.2026): antes decia "con las demas tablas por modelo de la medicion mas estricta"; de
% esas tablas quedan al final del anexo solo la del ciclo minimo y la de IPPO con mensajes

\subsection*{Comprobaciones}

% Integracion (23.10.2026): capitulo4.tex l. 1462-1482 del commit 7181ab8 (4.7.5), sin cambios. El parrafo de la
% auditoria, que sigue, remite a este ("mas arriba, en este apartado") en vez de repetir las 117 re-evaluaciones.
Re-evaluar una política o el actuado no cambia lo que ya estaba
medido. Con las mismas semillas, todas las columnas que ya existían
(\texttt{timeLoss}, espera, colas, paradas, \textit{throughput},
emisiones y métricas de corredor y de \textit{spillback}) se
reproducen dígito a dígito, y las columnas nuevas se agregan al final.
Esta comprobación se hizo al cerrar la medición más estricta, con las
117 evaluaciones que tenían una versión anterior: 110 de políticas, las
cinco del actuado y las del fijo de P1 y P2, cuyo programa no cambió.
Ninguna tiene un valor distinto en las 30 semillas. Solo el fijo de P3
a P5, que cambió de programa, da otros números; el del plan se conserva
como un control propio (Sección~\ref{sec:ejecucion}). Por eso, en P1 a P5, las diferencias
entre el análisis de robustez y el resultado principal salen solo de
la métrica y del fijo ajustado desde 10~s. En P6, cuyo fijo es el mismo en los
dos análisis, salen solo de la métrica.
% fuente: comparacion de cada plan2/parciales/eval_<esc>_<brazo>.csv re-evaluado contra su version en git
% (HEAD), revisada el 27.09.2026 despues del cierre: 120 archivos con version anterior, 117 con 0 celdas
% distintas en las columnas comunes (110 politicas, 5 actuados, fijo de corredor y corredor_alta) y 3
% distintos (fijo de red, red_alta y malla3); sin version anterior: el trivial de cada caso y las siete
% evaluaciones de spill que no se habian lanzado; nuevas: retraso_total_prom, retraso_insercion_prom y
% no_insertados (evaluar.METRICAS)

\noindent\textbf{Lo que dice la auditoría.} La medición más estricta
se auditó con tres preguntas. La primera, si cambió algún modelo: no.
Al volver a evaluar las políticas y el actuado con las mismas semillas,
todas las columnas que ya se medían se reproducen dígito a dígito
(más arriba, en este apartado), y el \texttt{timeLoss} de las 39
versiones evaluadas y del actuado de P1 a P5 es el del veredicto
pre-registrado. Solo el fijo de P3 a P5, que cambió de programa, da
otros números. La segunda, si algún error de simulación explica los
cambios: en las re-simulaciones de control, con los dos fijos y con
IPPO con mensajes, no hubo teletransportes, colisiones ni frenadas de
emergencia, y el \texttt{timeLoss} recalculado del \texttt{tripinfo}
coincide con el del archivo de resultados. La tercera, si lo ya
entrenado recupera la ventaja con otra selección: no. En P2, elegir la
versión guardada de IPPO con mensajes por el retraso total en la
validación da la misma que ya se había elegido, y ninguna de sus 60
versiones guardadas, tres entrenamientos por veinte validaciones, baja
del fijo en retraso total: la mejor queda un $0.98$\,\% por encima. En P3 a P5,
ninguna validación de ningún aprendiz queda por debajo del ciclo mínimo.
% fuente: Cap. IV, sec:remedicion (117 evaluaciones con version anterior, 0 celdas distintas); timeloss_prom de las
% politicas publicadas y del actuado frente a plan2/veredicto_congelado_0355.csv (diferencia maxima 0.0005);
% auditoria del 27.09.2026: seguridad.txt (14 re-simulaciones de P1 a P5 con tripinfo y estadisticas de SUMO:
% 0 teleports, 0 colisiones, 0 paradas y frenadas de emergencia; timeLoss del tripinfo igual al del CSV),
% analisis_p2.txt (validacion en 999 y 1999 por retraso total de los 60 checkpoints de ippo_vecinos en
% corredor_alta: base publicada s2042, ep120 por las dos metricas; minimo +0.98 % frente al fijo, s4042 ep90) y
% techo_validacion.txt (red, red_alta y malla3: 0 validaciones por debajo del trivial en todos los aprendices)

\subsection*{Qué queda en pie}

La ventaja frente al fijo
de los modelos con el actor de IPPO y mensajes es robusta en P1: con
las dos métricas y los dos fijos, IPPO con mensajes, MAPPO e IPPO con
penalización de \textit{spillback} funcionan, entre $-5.1$ y
$-7.1$\,\%. En P2 a P5 depende de la métrica y de la referencia. En P2
la ganancia en \texttt{timeLoss} va acompañada de más espera para
entrar a la red, y con el retraso total desaparece. En P3 y P4, un fijo
de ciclo corto que la rejilla del primer barrido no incluía iguala o
supera a las versiones evaluadas. En P5, el empate y la falla los
separa el tratamiento de las diferencias nulas en la prueba. En P6 la
medición no cambia ningún veredicto, porque el fijo es el mismo y la
espera para entrar es común al fijo, al actuado y a las versiones
evaluadas: con las dos métricas, ningún aprendiz del plan mejora al
fijo. Frente al actuado, en cambio, la ventaja de IPPO con mensajes se
mantiene en los seis casos y en el resumen en un número de P1 a P6: $-14.35$\,\% con el \texttt{timeLoss} y
$-11.45$\,\% con el retraso total ($-7.72$ y $-8.71$\,\% sobre P1 a
P5). Con lo ya entrenado la ventaja frente al fijo no se recupera, y
no está medido si un reentrenamiento con el retraso total en la
validación la recuperaría. El reentrenamiento del
Anexo~G corrige otra cosa, el descuento y la
exploración, y se mide con otras semillas. El veredicto del Capítulo~\ref{cap:resultados} sigue siendo el
pre-registrado; este anexo acota hasta dónde llega.
% fuente: plan2/veredicto_congelado_0355.csv y plan2/robustez_veredicto.csv (P1: ippo_vecinos, mappo e ippo_spill
% entre -5.1 y -7.1 en las dos; P6: mismos veredictos que plan2/veredicto.csv); plan2/sintesis_numero_unico.csv
% (ippo_vecinos frente al actuado -14.35, sin_P6 -7.72) y plan2/robustez_sintesis_numero_unico.csv (-11.45, sin_P6
% -8.71)

\subsection*{Barrido de demanda en el corredor, sin reentrenar}

El barrido de demanda de la Sección~\ref{sec:analisis-complementarios}
evalúa, en la red de P1 y con la demanda multiplicada por factores de
0.5 a 2.25 o con un perfil variable, las versiones evaluadas de IPPO
con mensajes de P1 y de P2, sin reentrenarlas, el actuado y el plan
fijo de P1 sin volver a ajustar. La referencia de cada nivel es un fijo
ajustado desde 10~s para ese nivel, con el mismo procedimiento que en
el resto de este anexo: una rejilla que parte del verde mínimo de 10~s,
las semillas 2001 a 2003 y el retraso total como criterio. Con
$\times$1.0 y $\times$2.0 ese barrido elige los mismos programas que
los fijos ajustados de P1 y P2,
así que en esos niveles el cambio en \texttt{timeLoss} coincide con el
de la Tabla~\ref{tab:plan2-resumen}: $-6.8$ y $-23.3$\,\% para las dos
políticas. La Tabla~\ref{tab:plan2-demanda} y la
Figura~\ref{fig:plan2-demanda} dan el cambio en las dos métricas. Su
resumen está en la Sección~\ref{sec:sensibilidad-demanda}.
% fuente: plan2/barrido_corredor_f<nivel>.csv (ganador por retraso_total_prom; malla con verde de avenida desde
% 10 s; f100: 20/10 0/15 y f200: 25/15 0/25, los programas del fijo de P1 y P2 del 25.09);
% plan2/demanda/resumen_demanda.csv (delta_timeloss_prom de ippo_vecinos_s4042 en corredor_f100, -6.82, y de
% ippo_vecinos_s2042 en corredor_f200, -23.3)

% fuente: python barrido_demanda.py resumir (desde demo_rl/) -> plan2/demanda/resumen_demanda.csv
% Generado por barrido_demanda.py resumir. No editar a mano.
\begin{table}[H]
\setstretch{1}
\centering
\scriptsize
\setlength{\tabcolsep}{3pt}
\begin{tabular}{lrrrrrrrrr}
\toprule
 & & \multicolumn{4}{c}{$\Delta$ retraso total (\%)} & \multicolumn{4}{c}{$\Delta$ timeLoss (\%)} \\
\cmidrule(lr){3-6}\cmidrule(lr){7-10}
Demanda & \makecell{Fijo\\(s/veh)} & \makecell{Plan\\de P1} & \makecell{Actuado} & \makecell{IPPO\\P1} & \makecell{IPPO\\P2} & \makecell{Plan\\de P1} & \makecell{Actuado} & \makecell{IPPO\\P1} & \makecell{IPPO\\P2} \\
\midrule
$\times$0.5 & 15.9 & +5.9 & +1.5 & +3.9 & +56.7 & +6.0 & +1.5 & +3.9 & +58.1 \\
$\times$0.75 & 18.7 & +3.9 & ($-$0.2) & +3.5 & +61.0 & +4.0 & ($-$0.2) & +3.6 & +63.0 \\
$\times$1.0 & 21.7 & (+0.0) & ($-$1.0) & $-$6.7 & +70.3 & (+0.0) & ($-$1.1) & $-$6.8 & +71.5 \\
$\times$1.25 & 29.1 & +100.6 & +81.0 & +86.2 & +160.9 & +93.8 & +79.3 & +71.7 & +146.1 \\
$\times$1.5 & 57.5 & +211.9 & +163.3 & +121.2 & +149.9 & +122.1 & +45.0 & +61.9 & +64.3 \\
$\times$1.75 & 259.6 & +43.2 & +4.3 & +18.1 & $-$3.6 & +34.5 & $-$24.6 & +9.0 & $-$22.8 \\
$\times$2.0 & 394.3 & +50.7 & +13.0 & +53.8 & +1.6 & +53.8 & $-$19.8 & +34.8 & $-$23.3 \\
$\times$2.25 & 552.2 & +51.6 & +15.4 & +81.0 & +10.0 & +62.6 & $-$12.2 & +66.9 & $-$17.3 \\
variable & 180.2 & +46.4 & +28.9 & +21.8 & +20.6 & +27.7 & $-$18.9 & ($-$1.7) & $-$16.6 \\
\bottomrule
\end{tabular}
\caption[Barrido de demanda en el corredor]{Barrido de demanda en la red de P1: cambio del retraso de cada control frente al fijo ajustado desde 10~s para ese nivel (semillas de evaluación 1001 a 1030, IC95\,\% bootstrap pareado). Entre paréntesis, los cambios que no son significativos: su $p$ de Wilcoxon pareado es de 0.05 o más. Fijo: retraso total del fijo ajustado desde 10~s. Plan de P1: el programa fijo de P1 sin volver a ajustar. IPPO P1 y P2: las versiones evaluadas de IPPO con mensajes entrenadas con $\times$1.0 y $\times$2.0, sin reentrenar. Variable: bloques de 10~min con $\times$1.0, 1.5, 2.0, 2.0, 1.5 y 1.0.}
\label{tab:plan2-demanda}
\end{table}


% 23.10.2026: flota ([!tbp]); en [H] dejaba en blanco media pagina
\begin{figure}[!tbp]
\centering
\includegraphics[width=\textwidth]{fig_plan2_demanda.png}
\caption[Barrido de demanda en el corredor: cambio por control]{Barrido de demanda en la red de P1: cambio porcentual de cada control frente al fijo ajustado desde 10~s para ese nivel de demanda, en retraso total (panel superior) y en \texttt{timeLoss} (panel inferior), sobre las semillas de evaluación 1001 a 1030. Filas: el plan fijo de P1 sin volver a ajustar, el actuado y las versiones evaluadas de IPPO con mensajes entrenadas con $\times$1.0 (P1) y con $\times$2.0 (P2), sin reentrenar. Columnas: factor que multiplica toda la demanda de P1, y el perfil variable. Color: rojo, más retraso que ese fijo; azul, menos; la intensidad crece con el tamaño del cambio.}
\label{fig:plan2-demanda}
\end{figure}
% fuente: figura escrita por barrido_demanda.py resumir desde plan2/demanda/resumen_demanda.csv
% Integracion de las tareas 4 y 5 (23.10.2026): barrido_demanda.py se paso a los nombres del glosario y la figura se
% regenero (rotulos "Fijo de P1 sin reajustar" e "IPPO con mensajes entrenado con x1.0 (P1)" / "x2.0 (P2)"; titulos
% "frente al fijo ajustado desde 10 s"; mismos datos). Por eso se quito la ultima frase del caption, que decia: "La
% figura usa los nombres anteriores: fijo re-barrido es el fijo ajustado desde 10 s; sin retunear, sin volver a
% ajustar, e IPPO vecinos, IPPO con mensajes."

La política de P1 solo gana con la demanda con que se entrenó. En
\texttt{timeLoss}, con $\times$1.0 reduce el retraso un $6.8$\,\%, y
con $\times$0.5 y $\times$0.75 pierde un $3.9$ y un $3.6$\,\%. Con
esas demandas el fijo ajustado desde 10~s es un ciclo corto sin desfases
(10/10~s con $\times$0.5 y 12/10~s con $\times$0.75), y la política no
lo alcanza, aunque queda por delante del plan de P1 sin volver a ajustar, que
pierde un $6.0$ y un $4.0$\,\%. La política de P2 gana en
\texttt{timeLoss} con $\times$1.75, $\times$2.0 y $\times$2.25 y con el
perfil variable, entre un $16.6$ y un $23.3$\,\%, y con $\times$0.5 a
$\times$1.0 pierde entre un $58.1$ y un $71.5$\,\%: una política
entrenada con la avenida saturada no sirve con poca demanda.
% fuente: plan2/demanda/resumen_demanda.csv (delta_timeloss_prom de ippo_vecinos_s4042, ippo_vecinos_s2042 y
% fijo_p1); plan2/barrido_corredor_f050.csv (ganador 10/10 0/0) y plan2/barrido_corredor_f075.csv (12/10 0/0)

En retraso total, que suma la espera para entrar a la red, la ganancia
con demanda alta baja mucho o se vuelve pérdida. Con $\times$1.75 a $\times$2.25, la
política de P2 pasa a $-3.6$, $+1.6$ y $+10.0$\,\%, y el actuado, que
en \texttt{timeLoss} ganaba entre un $12.2$ y un $24.6$\,\%, pierde
entre un $4.3$ y un $15.4$\,\%. Con el perfil variable, que pasa de
$\times$1.0 a $\times$2.0 y vuelve dentro de la hora, pierden todos en
retraso total, entre un $20.6$ y un $46.4$\,\%. Es la misma cola que
queda fuera de la red en P2 (más arriba, en ``Lectura por caso''): con
demanda
por encima de la capacidad, la ventaja en \texttt{timeLoss} de los
controles adaptativos no se sostiene cuando se cuenta la espera para
entrar. Con $\times$1.0 o menos la espera para entrar es pequeña y las
dos métricas dan casi lo mismo.
% fuente: plan2/demanda/resumen_demanda.csv (delta_retraso_total_prom y delta_timeloss_prom de actuado e
% ippo_vecinos_s2042 en corredor_f175, f200 y f225; corredor_var: delta_retraso_total_prom de +20.62 a +46.41).
% Verificacion de las tareas 4 y 5 (23.10.2026): decia "la ganancia con demanda alta casi desaparece", pero con x2.0
% y x2.25 la politica de P2 pierde de forma significativa (ippo_vecinos_s2042, delta_retraso_total_prom +1.63 [0.64,
% 2.56], p 0.00219, y +10.04 [8.61, 11.43], p 1.86e-09); con x1.75 gana un 3.64 % (p 3.24e-06)

Entre $\times$1.25 y $\times$1.5 colapsan todos los controles salvo el
fijo ajustado desde 10~s, que en esos niveles pasa a 25/10~s, con desfase de
25~s en $\times$1.25 y sin desfase en $\times$1.5: frente a él pierden
entre un $45.0$ y un $146.1$\,\% en \texttt{timeLoss}. La explicación
más probable del colapso es el giro a la izquierda permitido en el
carril único de la avenida, que el programa da en verde sin prioridad
(\texttt{g}, Sección~\ref{sec:construccion}): un vehículo que espera
un hueco en el tránsito opuesto para girar detiene a los que vienen
detrás en el mismo carril. En una comprobación de la fase de
diagnóstico, con diez semillas y fuera del protocolo del plan, el
colapso desaparecía al convertir esos giros en movimientos rectos. Esa
comprobación no forma parte de los resultados registrados del plan, y
la explicación queda como probable. El barrido mide hasta dónde se
sostiene una política fuera de su demanda de entrenamiento, y usa un
único escenario base (Sección~\ref{sec:analisis-complementarios}). No
mide si el método funciona con otra demanda, lo que exige entrenar en
cada nivel, como en los niveles de la Tabla~\ref{tab:efecto-demanda}.
% fuente: plan2/demanda/resumen_demanda.csv (corredor_f125 y corredor_f150, delta_timeloss_prom: fijo_p1 +93.79 y
% +122.13, actuado +79.34 y +45.01, ippo_vecinos_s4042 +71.65 y +61.88, ippo_vecinos_s2042 +146.09 y +64.27);
% plan2/barrido_corredor_f125.csv (25/10 0/25) y plan2/barrido_corredor_f150.csv (25/10 0/0). La comprobacion
% con los giros convertidos en rectos (diagnostico del 27.09.2026) no tiene archivo en el repositorio: para
% citarla como resultado hay que repetirla con el protocolo del plan. Verificacion de las tareas 4 y 5 (23.10.2026):
% decia "lo que exigiria entrenar en cada nivel" sin decir que eso ya se hizo; ahora remite a los niveles con el
% modelo entrenado en cada uno (capitulo5.tex, sec:efecto-demanda, Tabla tab:efecto-demanda)

\subsection*{Tablas por modelo}

% Integracion (23.10.2026): bloque "Tablas del analisis de robustez" del Anexo A (anexos.tex l. 134-197 del commit
% 7181ab8), con sus 11 inclusiones, que aqui pasan a F.6 en adelante y salen del Anexo A (cada tabla generada se
% incluye una sola vez). De su primer parrafo (l. 136-152) queda solo la primera frase, con "las anteriores" cambiado
% por "las del Anexo A" y la remision a sec:robustez cambiada por "mas arriba, en Que cambia", porque el resto ya esta
% en la entrada y en "Que cambia". El segundo parrafo (l. 160-186) no se pega, porque repite "Lectura por caso"; su
% unica cifra que no esta en el texto de este anexo, el retraso de insercion del trivial en P6 (99 s/veh), esta en la
% fuente del parrafo de P6 de "Lectura por caso" (99.185 en plan2/evaluacion_malla3_alta_resumen.csv).
% Reduccion del Anexo F (23.10.2026; plan del 14.10, aprobado por Fabian: de la robustez por modelo queda solo la
% tabla de IPPO con mensajes). Este bloque tenia 18 tablas (F.6 a F.23, unas 18 paginas con 25-45 % en blanco): las
% 10 por modelo y las 8 de entrenamientos. Quedan tres: (1) la del ciclo minimo, que se usa en la lectura de P3 y en
% el parrafo del ciclo minimo de "Resumen en un numero y reparto de la variacion"; (2) la de IPPO con mensajes; (3) la
% de entrenamientos de Double DQN, solo porque recomendaciones.tex (l. 72) cita tab:robustez-veredicto-bases-dqn y
% ese archivo no es de esta tarea. Si esa remision pasa a texto, la tercera se puede quitar. Salen las tablas por
% modelo del actuado, Q-learning, SARSA, Q-learning (colas), Double DQN, IPPO sin mensajes, IPPO con penalizacion de
% spillback y MAPPO, y las otras siete de entrenamientos. Ninguna de sus etiquetas se citaba fuera de este anexo
% (grep de tab:robustez-veredicto-* en TESIS/**/*.tex, 23.10.2026); la frase que las presentaba ("Las Tablas
% tab:robustez-veredicto-actuado a tab:robustez-veredicto-bases-mappo repiten las del Anexo A en P1 a P6 con la
% medicion del analisis de robustez") pasa a la que sigue. Los .tex generados siguen en tablas/ y "python
% veredicto.py --prefijo robustez_" los sigue escribiendo; el texto remite a robustez_veredicto_bases.tex, asi que
% ese archivo no se debe borrar. Verificacion de las tareas 4 y 5 (23.10.2026): sale el inciso "a la que remiten las
% recomendaciones" de la frase de la tabla de DQN, que era una nota de edicion; el motivo queda en este comentario.
\begin{sloppypar}
La Tabla~\ref{tab:robustez-veredicto-trivial} da el veredicto del
ciclo mínimo y la Tabla~\ref{tab:robustez-veredicto-ippo-vecinos}, el
de IPPO con mensajes, en P1 a P6 y con la medición del análisis de
robustez (más arriba, en ``Qué cambia''). La
Tabla~\ref{tab:robustez-veredicto-bases-dqn} da los tres entrenamientos
de Double DQN en cada caso. Las demás tablas por modelo y de
entrenamientos no se incluyen. El
veredicto de cada modelo en cada caso, con sus intervalos, sus valores~$p$,
el cambio frente al actuado y cuántos de los tres entrenamientos
convergen, está en el repositorio del proyecto (Anexo~B), en
\texttt{plan2/\allowbreak robustez\_veredicto.csv} y
\texttt{plan2/\allowbreak robustez\_veredicto.md}, y las tablas de
entrenamientos de todos los modelos, en
\texttt{TESIS/\allowbreak secciones/\allowbreak tablas/\allowbreak robustez\_veredicto\_bases.tex}.
\end{sloppypar}
% fuente: hechos_v21.md; plan2/robustez_veredicto.csv (28.09.2026 09:15, P1-P6; en P1-P5 igual a
% plan2/veredicto_v21_respaldo.csv); P6: hechos_2509.md (fijo_2509 = fijo re-barrido = 36.726 s/veh de
% timeLoss) y cabecera de plan2/malla3_alta.add.xml (avenida 10-35 s, transversal 10-25 s);
% plan2/barrido_malla3_alta.csv (72 filas; el menor timeloss_prom, 36.91, y el menor retraso_total_prom,
% 82.63, son la misma fila: avenida 15 s, transversal 10 s, desfases cero);
% se regeneran con "python veredicto.py --prefijo robustez_" desde demo_rl/
% fuente de lo que hay en el repositorio: plan2/robustez_veredicto.csv (una fila por caso y modelo, 60 filas, con
% delta_fijo, ic_fijo, p_fijo, delta_act, ic_act, p_act, convergio y convergen, k de 3) y plan2/robustez_veredicto.md
% (la matriz y una tabla por modelo: actuado, ciclo minimo y los ocho aprendices); los dos los escribe veredicto.py
% --prefijo robustez_ (28.09.2026; nombres internos sin cambios). Las tablas de entrenamientos (validacion elegida,
% retraso y cambio de cada base) no estan en esos dos archivos: estan en tablas/robustez_veredicto_bases.tex
% (veredicto.py; datos de plan2/evaluacion_<esc>_resumen.csv y de los meta de cada base)
% Generado por veredicto.py a partir de plan2/. No editar a mano.
\begin{table}[H]
\setstretch{1}
\centering
\scriptsize
\setlength{\tabcolsep}{1.5pt}
\begin{tabular}{llcrrrrrrcl}
\toprule
Caso & Escenario & Sem. & Carr. & \makecell{Dem.\\(veh/h)} & \makecell{Retraso\\(s/veh)} & \makecell{$\Delta$\\retraso\\vs fijo} & \makecell{$\Delta$\\retraso\\vs actuado} & \makecell{$\Delta$\\recompensa} & Conv. & Veredicto \\
\midrule
P1 & corredor & 2 & 1 & 1810 & \makecell[r]{106.26\\{}$\pm$ 2.08} & \makecell[r]{+390.3\,\%\\{}[379.8,\,400.7]\\$p$ $<$\,0.001} & -- & -- & -- & FALLA \\
P2 & corredor\_alta & 2 & 1 & 3620 & \makecell[r]{1055.41\\{}$\pm$ 11.71} & \makecell[r]{+167.7\,\%\\{}[163.7,\,171.6]\\$p$ $<$\,0.001} & -- & -- & -- & FALLA \\
P3 & red & \makecell[c]{3\\{}+ rot.} & 2 & 2150 & \makecell[r]{20.17\\{}$\pm$ 0.11} & \makecell[r]{$-$0.5\,\%\\{}[$-$0.8,\,$-$0.1]\\$p$ = 0.0128} & -- & -- & -- & MEJORA \\
P4 & red\_alta & \makecell[c]{3\\{}+ rot.} & 2 & 2580 & \makecell[r]{28.29\\{}$\pm$ 0.92} & \makecell[r]{$-$1.0\,\%\\{}[$-$2.4,\,0.5]\\$p$ = 0.217} & -- & -- & -- & EMPATA \\
P5 & malla3 & 9 & 2 & 4350 & \makecell[r]{21.25\\{}$\pm$ 0.05} & \makecell[r]{+0.0\,\%\\{}[0.0,\,0.0]\\$p$ = 1} & -- & -- & -- & EMPATA \\
P6 & malla3\_alta & 9 & 2 & 8700 & \makecell[r]{276.72\\{}$\pm$ 46.15} & \makecell[r]{+236.7\,\%\\{}[186.5,\,290.3]\\$p$ $<$\,0.001} & -- & -- & -- & FALLA \\
\bottomrule
\end{tabular}
\caption[Robustez: veredicto de Trivial]{Robustez. Trivial: veredicto por caso. Sem.: semáforos (rot.: más una rotonda); Carr.: carriles por sentido de la avenida; Dem.: demanda total. Retraso: retraso total (timeLoss más la espera para entrar a la red), media $\pm$ IC95\,\% sobre las semillas de evaluación 1001 a 1030; $\Delta$ con IC95\,\% bootstrap pareado y $p$ de Wilcoxon pareado. Control de cordura: el programa fijo con cada verde en el mínimo del agente (10~s) y los desfases en 0, es decir, cambiar de fase en cuanto la máscara lo permite. Fila de referencia juzgada frente al fijo tuneado.}
\label{tab:robustez-veredicto-trivial}
\end{table}

% Generado por veredicto.py a partir de plan2/. No editar a mano.
\begin{table}[H]
\setstretch{1}
\centering
\scriptsize
\setlength{\tabcolsep}{1.5pt}
\begin{tabular}{llcrrrrrrcl}
\toprule
Caso & \makecell[l]{Red y\\demanda} & Sem. & Carr. & \makecell{Dem.\\(veh/h)} & \makecell{Retraso\\(s/veh)} & \makecell{Cambio\\frente al\\fijo} & \makecell{Cambio\\frente al\\actuado} & \makecell{Cambio de\\la recom-\\pensa} & Conv. & Veredicto \\
\midrule
P1 & \makecell[l]{Corredor,\\{}nominal} & 2 & 1 & 1810 & \makecell[r]{20.22\\{}$\pm$ 0.17} & \makecell[r]{$-$6.7\,\%\\{}[$-$7.8,\,$-$5.7]\\$p$ $<$\,0.001} & \makecell[r]{$-$5.7\,\%\\{}[$-$6.6,\,$-$4.8]\\$p$ $<$\,0.001} & \makecell[r]{+7.0\,\%\\{}[5.9,\,8.3]\\$p$ $<$\,0.001} & \makecell[c]{sí\\{}(3/3)} & FUNCIONA \\
P2 & \makecell[l]{Corredor,\\{}alta} & 2 & 1 & 3620 & \makecell[r]{400.68\\{}$\pm$ 3.64} & \makecell[r]{+1.6\,\%\\{}[0.6,\,2.6]\\$p$ = 0.00219} & \makecell[r]{$-$10.1\,\%\\{}[$-$11.2,\,$-$8.8]\\$p$ $<$\,0.001} & \makecell[r]{+19.9\,\%\\{}[19.1,\,20.7]\\$p$ $<$\,0.001} & \makecell[c]{sí\\{}(2/3)} & FALLA \\
P3 & \makecell[l]{Red con\\{}rotonda,\\{}nominal} & \makecell[c]{3\\{}+ rot.} & 2 & 2150 & \makecell[r]{20.17\\{}$\pm$ 0.11} & \makecell[r]{$-$0.5\,\%\\{}[$-$0.8,\,$-$0.1]\\$p$ = 0.0128} & \makecell[r]{$-$7.6\,\%\\{}[$-$8.2,\,$-$7.0]\\$p$ $<$\,0.001} & \makecell[r]{+0.6\,\%\\{}[0.3,\,0.9]\\$p$ $<$\,0.001} & \makecell[c]{sí\\{}(3/3)} & \makecell[l]{MEJORA\\{}(igual al\\{}ciclo mínimo)} \\
P4 & \makecell[l]{Red con\\{}rotonda,\\{}alta} & \makecell[c]{3\\{}+ rot.} & 2 & 2580 & \makecell[r]{29.30\\{}$\pm$ 0.92} & \makecell[r]{+2.6\,\%\\{}[0.8,\,4.4]\\$p$ = 0.00984} & \makecell[r]{$-$3.5\,\%\\{}[$-$4.7,\,$-$2.3]\\$p$ $<$\,0.001} & \makecell[r]{$-$8.2\,\%\\{}[$-$9.0,\,$-$7.4]\\$p$ $<$\,0.001} & \makecell[c]{sí\\{}(3/3)} & FALLA \\
P5 & \makecell[l]{Malla $3 \times 3$,\\{}nominal} & 9 & 2 & 4350 & \makecell[r]{21.68\\{}$\pm$ 0.35} & \makecell[r]{+2.0\,\%\\{}[0.7,\,3.6]\\$p$ = 0.0277} & \makecell[r]{$-$16.1\,\%\\{}[$-$17.3,\,$-$14.7]\\$p$ $<$\,0.001} & \makecell[r]{$-$2.8\,\%\\{}[$-$5.0,\,$-$0.9]\\$p$ = 0.0252} & \makecell[c]{sí\\{}(3/3)} & FALLA \\
P6 & \makecell[l]{Malla $3 \times 3$,\\{}alta} & 9 & 2 & 8700 & \makecell[r]{83.05\\{}$\pm$ 1.47} & \makecell[r]{+1.1\,\%\\{}[$-$0.6,\,2.8]\\$p$ = 0.596} & \makecell[r]{$-$23.9\,\%\\{}[$-$25.4,\,$-$22.3]\\$p$ $<$\,0.001} & \makecell[r]{$-$7.8\,\%\\{}[$-$12.2,\,$-$3.9]\\$p$ = 0.00322} & \makecell[c]{no\\{}(1/3)} & \makecell[l]{PENDIENTE\\{}(no convergió;\\{}provisional:\\{}EMPATA)} \\
\bottomrule
\end{tabular}
\caption[Robustez: veredicto de IPPO con mensajes]{Robustez. IPPO con mensajes: veredicto por caso. Sem.: semáforos (rot.: más una rotonda); Carr.: carriles por sentido de la avenida; Dem.: demanda total. Retraso: retraso total (timeLoss más la espera para entrar a la red), media $\pm$ IC95\,\% sobre las semillas de evaluación 1001 a 1030; cambio (\%) con IC95\,\% bootstrap pareado y $p$ de Wilcoxon pareado. Versión evaluada: la versión guardada del entrenamiento mediano de los tres por la validación (semillas 999 y 1999). Cambio de la recompensa: el de la recompensa propia del modelo (ret\_retraso) frente a la del fijo, positivo = mejor. Conv.: si convergió el entrenamiento de la versión evaluada; entre paréntesis, cuántos de los tres entrenamientos convergen. Igual al ciclo mínimo: la política da el mismo retraso y la misma espera que el ciclo mínimo.}
\label{tab:robustez-veredicto-ippo-vecinos}
\end{table}


% Tabla de entrenamientos de Double DQN: copia literal del bloque de DQN de tablas/robustez_veredicto_bases.tex (l.
% 111-148 del archivo escrito por "python veredicto.py --prefijo robustez_"; sin cambios desde el commit b9c4e91). Ese
% archivo trae las ocho tablas de entrenamientos en una sola inclusion, y esta es la unica que se cita fuera del anexo
% (recomendaciones.tex, l. 72), por eso se copia sola. Si veredicto.py vuelve a escribir ese archivo, hay que copiar
% de nuevo el bloque; si la remision de recomendaciones.tex pasa a texto, esta tabla se puede quitar.
% fuente: plan2/evaluacion_<esc>_resumen.csv (filas dqn_s42, dqn_s2042 y dqn_s4042, metrica retraso_total_prom:
% media, ic95, delta_pct, delta_ic95 y p_wilcoxon) y plan2/politica_<esc>_dqn_s<base>.json (val_elegido y
% ep_elegido; convergencia: veredicto.convergencia sobre val_serie). P2 revisado el 23.10.2026: retraso 411.456,
% 388.6143 y 461.124 s/veh (rango 72.51) y validacion 123.24 (ep 190), 121.655 (ep 200) y 122.485 (ep 50)
\begin{table}[H]
\setstretch{1}
\centering
\scriptsize
\setlength{\tabcolsep}{3pt}
\begin{tabular}{lllrlrrr}
\toprule
Caso & Red y demanda & \makecell{Semilla del\\entrenamiento} & \makecell{Validación\\elegida\\(s/veh)} & \makecell{Conver-\\gió} & \makecell{Retraso\\evaluación\\(s/veh)} & \makecell{Cambio frente\\al fijo [IC95\,\%]} & $p$ \\
\midrule
P1 & \makecell[l]{Corredor,\\{}nominal} & 42 & \makecell[r]{20.54\\{}(ep 120)} & sí & \makecell[r]{21.16\\{}$\pm$ 0.12} & \makecell[r]{$-$2.4\,\%\\{}[$-$3.5,\,$-$1.2]} & 0.00124 \\
P1 & \makecell[l]{Corredor,\\{}nominal} & 2042 & \makecell[r]{20.13\\{}(ep 190)} & sí & \makecell[r]{20.67\\{}$\pm$ 0.13} & \makecell[r]{$-$4.6\,\%\\{}[$-$5.6,\,$-$3.6]} & $<$\,0.001 \\
P1 & \makecell[l]{Corredor,\\{}nominal} & \makecell[l]{4042\\{}(evaluada)} & \makecell[r]{20.32\\{}(ep 180)} & sí & \makecell[r]{20.94\\{}$\pm$ 0.16} & \makecell[r]{$-$3.4\,\%\\{}[$-$4.8,\,$-$2.2]} & $<$\,0.001 \\
 & & \makecell[l]{rango entre\\{}los tres} & & & 0.49 & & \\
P2 & \makecell[l]{Corredor,\\{}alta} & 42 & \makecell[r]{123.24\\{}(ep 190)} & sí & \makecell[r]{411.46\\{}$\pm$ 5.45} & \makecell[r]{+4.4\,\%\\{}[2.8,\,5.9]} & $<$\,0.001 \\
P2 & \makecell[l]{Corredor,\\{}alta} & 2042 & \makecell[r]{121.66\\{}(ep 200)} & sí & \makecell[r]{388.61\\{}$\pm$ 6.06} & \makecell[r]{$-$1.4\,\%\\{}[$-$2.9,\,0.1]} & 0.0732 \\
P2 & \makecell[l]{Corredor,\\{}alta} & \makecell[l]{4042\\{}(evaluada)} & \makecell[r]{122.48\\{}(ep 50)} & sí & \makecell[r]{461.12\\{}$\pm$ 5.70} & \makecell[r]{+17.0\,\%\\{}[15.4,\,18.5]} & $<$\,0.001 \\
 & & \makecell[l]{rango entre\\{}los tres} & & & 72.51 & & \\
P3 & \makecell[l]{Red con rotonda,\\{}nominal} & 42 & \makecell[r]{21.12\\{}(ep 100)} & sí & \makecell[r]{21.85\\{}$\pm$ 0.10} & \makecell[r]{+7.8\,\%\\{}[7.2,\,8.5]} & $<$\,0.001 \\
P3 & \makecell[l]{Red con rotonda,\\{}nominal} & 2042 & \makecell[r]{20.86\\{}(ep 190)} & sí & \makecell[r]{21.44\\{}$\pm$ 0.13} & \makecell[r]{+5.8\,\%\\{}[5.2,\,6.5]} & $<$\,0.001 \\
P3 & \makecell[l]{Red con rotonda,\\{}nominal} & \makecell[l]{4042\\{}(evaluada)} & \makecell[r]{20.91\\{}(ep 200)} & sí & \makecell[r]{21.68\\{}$\pm$ 0.10} & \makecell[r]{+7.0\,\%\\{}[6.4,\,7.6]} & $<$\,0.001 \\
 & & \makecell[l]{rango entre\\{}los tres} & & & 0.41 & & \\
P4 & \makecell[l]{Red con rotonda,\\{}alta} & 42 & \makecell[r]{28.18\\{}(ep 80)} & sí & \makecell[r]{29.33\\{}$\pm$ 1.02} & \makecell[r]{+2.7\,\%\\{}[0.9,\,4.5]} & 0.00761 \\
P4 & \makecell[l]{Red con rotonda,\\{}alta} & 2042 & \makecell[r]{28.78\\{}(ep 170)} & no & \makecell[r]{29.29\\{}$\pm$ 1.06} & \makecell[r]{+2.5\,\%\\{}[1.0,\,4.2]} & 0.00662 \\
P4 & \makecell[l]{Red con rotonda,\\{}alta} & \makecell[l]{4042\\{}(evaluada)} & \makecell[r]{28.43\\{}(ep 200)} & sí & \makecell[r]{29.37\\{}$\pm$ 1.02} & \makecell[r]{+2.8\,\%\\{}[1.1,\,4.6]} & 0.0019 \\
 & & \makecell[l]{rango entre\\{}los tres} & & & 0.08 & & \\
P5 & \makecell[l]{Malla $3 \times 3$,\\{}nominal} & 42 & \makecell[r]{23.20\\{}(ep 170)} & sí & \makecell[r]{23.42\\{}$\pm$ 0.11} & \makecell[r]{+10.2\,\%\\{}[9.7,\,10.8]} & $<$\,0.001 \\
P5 & \makecell[l]{Malla $3 \times 3$,\\{}nominal} & 2042 & \makecell[r]{22.89\\{}(ep 160)} & sí & \makecell[r]{23.03\\{}$\pm$ 0.13} & \makecell[r]{+8.4\,\%\\{}[7.8,\,9.1]} & $<$\,0.001 \\
P5 & \makecell[l]{Malla $3 \times 3$,\\{}nominal} & \makecell[l]{4042\\{}(evaluada)} & \makecell[r]{23.04\\{}(ep 190)} & no & \makecell[r]{23.62\\{}$\pm$ 0.09} & \makecell[r]{+11.2\,\%\\{}[10.7,\,11.6]} & $<$\,0.001 \\
 & & \makecell[l]{rango entre\\{}los tres} & & & 0.58 & & \\
P6 & \makecell[l]{Malla $3 \times 3$,\\{}alta} & 42 & \makecell[r]{37.85\\{}(ep 160)} & no & \makecell[r]{84.79\\{}$\pm$ 1.12} & \makecell[r]{+3.2\,\%\\{}[2.0,\,4.5]} & $<$\,0.001 \\
P6 & \makecell[l]{Malla $3 \times 3$,\\{}alta} & \makecell[l]{2042\\{}(evaluada)} & \makecell[r]{37.23\\{}(ep 130)} & no & \makecell[r]{86.33\\{}$\pm$ 1.18} & \makecell[r]{+5.0\,\%\\{}[3.7,\,6.5]} & $<$\,0.001 \\
P6 & \makecell[l]{Malla $3 \times 3$,\\{}alta} & 4042 & \makecell[r]{36.85\\{}(ep 160)} & sí & \makecell[r]{83.38\\{}$\pm$ 0.63} & \makecell[r]{+1.4\,\%\\{}[0.9,\,2.1]} & $<$\,0.001 \\
 & & \makecell[l]{rango entre\\{}los tres} & & & 2.94 & & \\
\bottomrule
\end{tabular}
\caption[Robustez: entrenamientos de DQN]{Robustez. DQN: los tres entrenamientos de cada caso. La validación elegida es la media del timeLoss en las semillas 999 y 1999 de la versión guardada que se elige en cada entrenamiento; el retraso y el cambio frente al fijo, sobre las semillas 1001 a 1030. El entrenamiento marcado (evaluada) es el mediano por validación y da la versión evaluada. El rango entre los tres se compara con el IC95\,\% de una sola política.}
\label{tab:robustez-veredicto-bases-dqn}
\end{table}
